\documentclass[11pt,a4paper]{article}
\usepackage[T1]{fontenc}
\usepackage[utf8]{inputenc}
\usepackage[french]{babel}
\usepackage[default]{sourcesanspro}
\usepackage[scaled=0.95]{inconsolata}
\usepackage{textcomp}
\usepackage[margin=2.1cm]{geometry}
\usepackage{booktabs}
\usepackage{colortbl}
\usepackage[most]{tcolorbox}
\usepackage{enumitem}
\usepackage{fancyhdr}
\usepackage{lastpage}
\usepackage{fontawesome5}
\usepackage{xcolor}
\usepackage{array}
\usepackage{ragged2e}
\usepackage{titlesec}
\usepackage{multicol}
\setlength{\headheight}{16pt}

\linespread{1.04}
\setlength{\emergencystretch}{2em}

% ----- Palette Horse Haven -----
\definecolor{navy}{HTML}{1A2744}
\definecolor{royalblue}{HTML}{2B5A83}
\definecolor{lightblue}{HTML}{F2F6FB}
\definecolor{gold}{HTML}{D4C5A9}
\definecolor{green}{HTML}{006B3F}
\definecolor{red}{HTML}{CC3333}
\definecolor{graytext}{HTML}{667788}

\newcolumntype{L}[1]{>{\RaggedRight\arraybackslash}p{#1}}

% ----- En-tête / pied de page -----
\pagestyle{fancy}
\fancyhf{}
\fancyhead[L]{\small\color{navy}\textbf{Horse Haven}}
\fancyhead[R]{\small\color{graytext}Étude de marché --- Maroc --- Août 2026}
\fancyfoot[C]{\footnotesize\color{graytext}Page \thepage\ sur \pageref{LastPage}}
\renewcommand{\headrulewidth}{0.8pt}
\renewcommand{\headrule}{\color{gold}\hrule width\headwidth height 0.8pt}

% ----- Titres -----
\titleformat{\section}
  {\Large\bfseries\color{navy}}
  {\thesection.}
  {0.6em}
  {}
  [\vspace{-0.3em}\textcolor{gold}{\rule{\textwidth}{1.6pt}}]
\titlespacing*{\section}{0pt}{5mm}{2.5mm}

\titleformat{\subsection}
  {\large\bfseries\color{royalblue}}
  {\thesubsection}
  {0.6em}
  {}
\titlespacing*{\subsection}{0pt}{4mm}{1.5mm}

% ----- Listes -----
\setlist[itemize]{leftmargin=5mm, itemsep=0.6mm, label=\textcolor{royalblue}{\small\raisebox{0.6pt}{\textbullet}}}
\setlist[enumerate]{leftmargin=6mm, itemsep=1mm, label=\textcolor{navy}{\bfseries\arabic*.}}

% ----- Tableaux -----
\renewcommand{\arraystretch}{1.25}
\arrayrulecolor{navy}
\rowcolors{2}{lightblue!45}{white}

% ----- Boîtes -----
\newtcolorbox{actionbox}{title={\faLightbulb\ Action recommandée}, colback=green!7, colframe=green, coltitle=white, colbacktitle=green, boxrule=0.9pt, arc=2mm, left=3mm, right=3mm, top=2mm, bottom=2mm, breakable, fonttitle=\small\bfseries}
\newtcolorbox{warnbox}{title={\faExclamationTriangle\ Risques à anticiper}, colback=red!6, colframe=red, coltitle=white, colbacktitle=red, boxrule=0.9pt, arc=2mm, left=3mm, right=3mm, top=2mm, bottom=2mm, breakable, fonttitle=\small\bfseries}
\newtcolorbox{keybox}{title={\faKey\ Point clé}, colback=lightblue, colframe=royalblue, coltitle=white, colbacktitle=royalblue, boxrule=0.9pt, arc=2mm, left=3mm, right=3mm, top=2mm, bottom=2mm, breakable, fonttitle=\small\bfseries}
\newtcolorbox{formulabox}{title={\faCalculator\ Mécanique fiscale}, colback=white, colframe=navy, coltitle=white, colbacktitle=navy, boxrule=0.9pt, arc=1.5mm, left=4mm, right=4mm, top=2mm, bottom=2mm, fonttitle=\small\bfseries}

% ----- En-tête de tableau -----
\newcommand{\thead}[1]{\textcolor{white}{\textbf{#1}}}

\begin{document}

% ================= EN-TÊTE =================
\noindent
\begin{tcolorbox}[colback=navy, colframe=navy, boxrule=0pt, arc=3mm, left=4mm, right=4mm, top=3mm, bottom=3mm]
\centering
{\color{white}\LARGE\faHorse\ \ \textsc{Horse Haven}}\\[2mm]
{\color{gold}\Large\bfseries Étude de marché --- Maroc}\\[1mm]
{\color{white!70!royalblue}\itshape Données réelles \& scénarios concrets}
\end{tcolorbox}

\vspace{2mm}
\noindent
\begin{tcolorbox}[colback=lightblue, colframe=royalblue, boxrule=0.6pt, arc=2mm, left=3mm, right=3mm, top=1.5mm, bottom=1.5mm]
{\small\color{graytext} Compilée à partir de sources publiques : SOREC, Le360, MAP, Oxford Business Group, UNESCO, FRMSE, CMI, ANRT, HCP, Research \& Markets, douane.gov.ma --- \textbf{Août 2026}}
\end{tcolorbox}

% ================= 1. SECTEUR ÉQUIN =================
\section{Le secteur équin marocain --- Chiffres clés}

{\small
\noindent
\begin{tabular}{L{4.6cm} L{7.4cm} L{3.5cm}}
\toprule
\rowcolor{navy}
\thead{Indicateur} & \thead{Valeur} & \thead{Source}\\
\midrule
\textbf{Cheptel équin national} & \textbf{$\sim$170\,000 chevaux} & SOREC / fiches projets\\
\textbf{Contribution au PIB} & \textbf{0,6\,\% ($\sim$6 milliards MAD/an)} & Min. Agriculture / Le360 (2024)\\
\textbf{Emplois directs \& indirects} & \textbf{30\,000+} & Le360 (2024)\\
\textbf{Naissances annuelles (5 races)} & \textbf{$\sim$4\,300} (+30\,\% depuis 2011) & SOREC / LesEco\\
\textbf{Cavaliers recensés} & \textbf{$\sim$20\,000} & Fiches projets\\
\textbf{Centres / clubs équestres} & \textbf{150+} sur tout le territoire & MesLoisirs.ma (2026)\\
\textbf{Courses hippiques / an} & \textbf{2\,400 à 3\,000} sur 7 hippodromes & SOREC / OBG\\
\textbf{Troupes de Tbourida} & \textbf{$\sim$1\,000} nationales ; 330+ au championnat & UNESCO / FRMSE\\
\textbf{Chevaux dédiés Tbourida} & \textbf{$\sim$5\,900} & FRMSE (2025)\\
\textbf{Manifestations équestres / an} & \textbf{2\,400+} & Min. Agriculture\\
\textbf{Races principales} & Barbe, Arabe-Barbe, Pur-sang arabe, PS anglais, Anglo-arabe & SOREC\\
\bottomrule
\end{tabular}}

\subsection{Infrastructure SOREC (organisme public, créé en 2003)}
\begin{itemize}
  \item 7 hippodromes, 5 haras nationaux (El Jadida, Bouznika, Meknès, Oujda, Marrakech)
  \item 43 stations de monte, 300 étalons, 13\,000 juments saillies/an
  \item Centre National d'Insémination Artificielle Équine de Bouznika --- 1\textsuperscript{er} d'Afrique du Nord, agréé UE
  \item 200 professionnels au service des éleveurs ; application d'identification «~Faras~»
\end{itemize}

% ================= 2. SCÉNARIOS RÉELS =================
\section{Scénarios réels --- Là où se trouvent les clients}

\subsection{\faFlag\ Salon du Cheval d'El Jadida (événement phare)}
\begin{itemize}
  \item \textbf{16\textsuperscript{e} édition, 30 sept. -- 5 oct. 2025 : $\sim$150\,000 visiteurs} (2024 : 160\,000+ ; 200\,000 attendus en 2025)
  \item $\sim$50 pays, 700+ cavaliers, 1\,100 chevaux, 600+ journalistes
  \item Sous le Haut Patronage de SM le Roi Mohammed VI ; ouvert par le Prince Héritier Moulay El Hassan
  \item Compétitions clés : Grand Prix SM le Roi (CSO), Grand Prix SM le Roi (Tbourida), Coupe des Champions Barbes, Arabe-Barbes, «~Title Show~» pur-sang arabes (nouveau Cat. A en 2025)
  \item Thème 2025 : «~Le bien-être du cheval~» --- le bien-être animal comme angle marketing
\end{itemize}

\begin{actionbox}
Exposer au Salon (village exposants) ou sponsoriser un espace ; 150\,000 clients potentiels réunis 6 jours en octobre.
\end{actionbox}

\subsection{\faUsers\ Tbourida / Fantasia (patrimoine immatériel UNESCO depuis 2021)}
\begin{itemize}
  \item $\sim$1\,000 troupes au niveau national ; 330+ au championnat ; 5\,900 chevaux dédiés
  \item Festival Fête du Trône 2025 à Rabat : \textbf{500 cavaliers, 50+ sorbas}, 4 jours complets
  \item Festival de Sidi Rahal (13\textsuperscript{e} éd., 2025) : 42 sorbas, \textbf{800+ chevaux/cavaliers, $\sim$30\,000 visiteurs}
  \item Festival d'Ain Chkef / Moulay Yacoub : 42 sorbas, 600 chevaux
  \item Trophée Hassan II de Tbourida (Rabat, juin) : 24 sorbas d'élite (18 séniors + 6 juniors)
\end{itemize}

\begin{actionbox}
Les sorbas ont besoin de : selles, étuis de fusil, harnachements traditionnels, couvertures, produits de pansage. Partenariats avec les mokaddems (chefs de troupe) pour commandes groupées.
\end{actionbox}

\subsection{\faTrophy\ Courses hippiques (7 hippodromes)}
\begin{itemize}
  \item 2\,400 à 3\,000 courses/an ; 2\,852 chevaux en course en 2016 (+31\,\% vs 2015) ; $\sim$12 partants/course
  \item $\rightarrow$ Propriétaires et entraîneurs achètent : couvertures d'écurie, bandages, compléments alimentaires, produits de pansage.
\end{itemize}

\subsection{\faGraduationCap\ Écoles \& clubs d'équitation (150+ centres)}
\begin{itemize}
  \item Prix réels (MesLoisirs.ma 2026) : cours 200--500 MAD/h ; abonnement mensuel 2\,000--4\,000 MAD
  \item Clubs : Royal Club Équestre Marrakech, Casa Anfa (fondé en 1912), Bouskoura, Mohammedia, Mamora, Rabat/Témara, Fès, Meknès, Agadir, Essaouira (Diabat), Taroudant, Ifrane
  \item Clubs de polo : Marrakech (Royal Polo Club, Jnan Amar), Rabat --- initiation 500--1\,500 MAD/séance, écuries de 40 chevaux
\end{itemize}

\begin{actionbox}
Approvisionner les clubs en sellerie, équipement cavalier (casques, bottes, chaps), matériel d'écurie. Les clubs = acheteurs récurrents en volume (B2B2C).
\end{actionbox}

\subsection{\faSuitcase\ Tourisme équestre (nouveau et en croissance)}
\begin{itemize}
  \item Ifrane : \textbf{181 km d'itinéraires balisés} inaugurés en mai 2026 (région Fès-Meknès + coopération Centre-Val de Loire, 1,4~M~MAD investis) --- 1\textsuperscript{er} itinéraire national de tourisme équestre
  \item Objectif : $\sim$230 km avec boucles ; guides en formation avec la FFE/FITE
  \item Prix : balade 1h 150--600 MAD ; trek Atlas 3--7 jours 2\,500--20\,000 MAD
\end{itemize}

\begin{actionbox}
Vendre du matériel de randonnée (chemises anti-mouches, guêtres de transport, selles de bât) + les opérateurs touristiques comme clients B2B.
\end{actionbox}

% ================= 3. E-COMMERCE =================
\section{La réalité du e-commerce au Maroc}

{\small
\noindent
\begin{tabular}{L{5.2cm} L{6.6cm} L{3.7cm}}
\toprule
\rowcolor{navy}
\thead{Indicateur} & \thead{Valeur} & \thead{Source}\\
\midrule
\textbf{Paiements e-commerce 2025} & \textbf{18 milliards MAD (+34\,\% vs 2024)} & CMI / Wafir (2026)\\
\textbf{Acheteurs en ligne 2025} & \textbf{6,8 millions (+278\,\% vs 2020)} = 24\,\% des adultes & Wafir\\
\textbf{Panier moyen} & \textbf{380 MAD} & Wafir\\
\textbf{Fréquence d'achat} & \textbf{4,2 fois/an/acheteur} & Wafir\\
\textbf{Paiements 2025} & Carte \textbf{47\,\%}, COD \textbf{38\,\%}, M-wallet \textbf{12\,\%}, virement 3\,\% & CMI\\
\textbf{Part du mobile} & \textbf{73\,\%} des transactions (vs 38\,\% en 2020) & Wafir\\
\textbf{Marché total 2025} & 3,17 Mds \$US ; prévision 3,51 Mds \$US en 2029 & Research \& Markets\\
\textbf{Croissance} & +25 à 30\,\%/an depuis 5 ans & amine.ma / CMI\\
\textbf{Pénétration smartphone} & \textbf{87\,\%} & ANRT 2025\\
\textbf{Social commerce} & \textbf{62\,\%} utilisent les boutiques IG/TikTok /mois & amine.ma 2026\\
\textbf{Top plateformes} & Jumia (28\,\%), Marjane (18\,\%), Glovo (12\,\%), JaJa (8\,\%), Carrefour (6\,\%) & Wafir\\
\bottomrule
\end{tabular}}

\subsection{\faShoppingCart\ Ce que cela implique pour Horse Haven}
\begin{itemize}
  \item \textbf{La livraison contre remboursement reste incontournable (38\,\%)} : proposer COD + paiement carte CMI
  \item \textbf{Mobile-first} : la boutique doit être optimisée mobile (73\,\% des transactions)
  \item \textbf{M-wallets en hausse (12\,\%)} : envisager l'intégration de portefeuilles mobiles
  \item \textbf{Le social commerce est massif (62\,\%)} : vendre via Instagram/TikTok est un comportement standard
  \item \textbf{Panier moyen 380 MAD} = les produits à 200--600 MAD se vendent le mieux ; les gros achats (selles 3\,000+ MAD) nécessitent des options de paiement en plusieurs fois
  \item Les niches verticales (comme l'équitation) sont à \textbf{faible densité concurrentielle} --- «~grand ouvert~» selon amine.ma
\end{itemize}

% ================= 4. IMPORTATION =================
\section{Coûts d'importation --- Scénario chiffré réel (douane)}

\begin{formulabox}
\begin{tabular}{@{}l@{\hspace{0.8em}}l@{}}
\texttt{Valeur CIF} & = FOB + Fret + Assurance\\
\texttt{Droit d'Importation (DI)} & = CIF $\times$ taux (0 à 30\,\%)\\
\texttt{TPF (Taxe Parafiscale)} & = CIF $\times$ 0,25\,\%\\
\texttt{Base TVA} & = CIF + DI + TPF\\
\texttt{TVA} & = Base $\times$ 20\,\% (taux normal)\\
\textbf{Total} & = DI + TPF + TVA\\
\end{tabular}
\end{formulabox}

\subsection{\faPercent\ Taux de droits par catégorie (PLF 2026, circulaire ADII 6702/210)}
{\small
\noindent
\begin{tabular}{L{12cm} L{3cm}}
\toprule
\rowcolor{navy}
\thead{Catégorie} & \thead{Taux DI}\\
\midrule
Matières premières & 0\,\%\\
Équipements industriels & 2,5\,\%\\
Produits semi-finis & 10\,\%\\
Produits finis courants & \textbf{25\,\%} (sellerie, vêtements, harnachements)\\
Produits fortement protégés & 30\,\% (textiles possibles)\\
\bottomrule
\end{tabular}}

\begin{keybox}
Les produits d'\textbf{origine UE avec certificat EUR.1 = 0\,\% de droits} (Accord d'Association Maroc-UE 2000). La plupart des marques équestres sont européennes $\rightarrow$ importer de France/Allemagne/Espagne/Italie avec EUR.1 pour supprimer le droit de douane.\\[1mm]
La TVA à l'import est \textbf{récupérable} si vous êtes assujetti (obligatoire $>$ 2~M~MAD de CA pour le commerce).\\[1mm]
Fret exemple : conteneur 20 pieds Chine$\rightarrow$Casablanca : 2\,000--6\,000 \$US ; fret aérien = 10--20\,\% du FOB.
\end{keybox}

\subsection{\faEuroSign\ Exemple chiffré --- import de 5\,000\,\texteuro{} ($\approx$53\,500 MAD) de sellerie depuis la France (avec EUR.1)}
{\small
\noindent
\begin{tabular}{L{9.5cm} L{5cm}}
\toprule
\rowcolor{navy}
\thead{Élément} & \thead{Montant (MAD)}\\
\midrule
FOB & 53\,500\\
Fret + assurance & 4\,500\\
CIF & 58\,000\\
DI (0\,\% avec EUR.1) & 0\\
TPF (0,25\,\%) & 145\\
Base TVA & 58\,145\\
TVA import (20\,\%) & 11\,629\\
\midrule
\textbf{Total droits et taxes} & \textbf{11\,774} ($\sim$17\,\% du coût débarqué total de 69\,774)\\
\bottomrule
\end{tabular}}

Sans EUR.1 (ex. origine Chine) : DI 25\,\% = 14\,500 MAD $\rightarrow$ total $\approx$ 29\,174 MAD ($\sim$50\,\% du CIF). \textbf{Privilégier impérativement les fournisseurs UE.}

% ================= 5. CONCURRENCE =================
\section{Concurrence réelle --- Qui vend de l'équipement équestre au Maroc}

{\small
\noindent
\begin{tabular}{L{4.0cm} L{2.5cm} L{9.0cm}}
\toprule
\rowcolor{navy}
\thead{Acteur} & \thead{Type} & \thead{Observations}\\
\midrule
\textbf{Decathlon.ma} & Grande distribution & 16 produits équestres ; tapis de selle 89--349 MAD ; selles promo 1\,999--3\,999 MAD ; uniquement entrée de gamme\\
\textbf{Maroc Marechalerie SARL} (marechalpro.com) & Spécialiste B2B & Casablanca ; fers, clous, outils de maréchalerie + sellerie limitée (tapis, longe) ; le spécialiste existant le plus proche\\
\textbf{Cavalor Maroc} (cavalormaroc.com) & Distributeur d'aliments & Marque belge distribuée au Maroc, ex. Juniorix 20~kg = 285 MAD --- un marché premium existe\\
\textbf{Fibersoil} (Casablanca) & Surfaces équestres & Sables de carrière/manèges pour clubs (B2B infrastructure, pas de retail)\\
\textbf{Revendeurs Facebook/Instagram} & Informel & Nombreux petits vendeurs via DM + COD --- canal dominant actuel, sans confiance ni garantie\\
\textbf{Boutiques en ligne UE} (equi-clic.com, cheval-shop.com, equiphorse.com) & Concurrents UE & Livraison au Maroc via fret privé : lent, cher, sans support local\\
\bottomrule
\end{tabular}}

\subsection{\faSearch\ Analyse des lacunes $\rightarrow$ positionnement Horse Haven}
\begin{enumerate}
  \item \textbf{Aucune boutique en ligne marocaine spécialisée} en équipement équestre n'existe (seulement maréchalerie + aliments)
  \item Decathlon couvre uniquement l'entrée de gamme ; les pros achètent à l'étranger ou via des revendeurs informels
  \item La demande B2B (150 clubs, 1\,000 sorbas, 7 hippodromes) n'est pas servie en ligne
  \item Origine UE + EUR.1 = import sans droits = prix compétitifs face aux revendeurs informels
  \item L'équitation est une \textbf{niche à faible densité} dans le marché e-commerce qui croît le plus vite d'Afrique
\end{enumerate}

% ================= 6. PRIX =================
\section{Repères de prix réels au Maroc (MAD)}
{\small
\noindent
\begin{tabular}{L{8.2cm} L{4.3cm} L{3cm}}
\toprule
\rowcolor{navy}
\thead{Article / Service} & \thead{Fourchette (MAD)} & \thead{Source}\\
\midrule
Tapis de selle (numnah) & 89--349 & Decathlon.ma\\
Selle entrée de gamme (promo) & 1\,999--3\,999 & Decathlon.ma\\
Aliment : Cavalor Juniorix 20~kg & 285 & cavalormaroc.com\\
Cours d'équitation 1h & 200--500 & MesLoisirs.ma\\
Abonnement mensuel club & 2\,000--4\,000 & MesLoisirs.ma\\
Balade 1h (tourisme) & 150--600 & MesLoisirs.ma\\
Randonnée week-end & 500--1\,500 & MesLoisirs.ma\\
Panier e-commerce moyen & 380 & Observatoire Wafir\\
\bottomrule
\end{tabular}}

% ================= 7. STRATÉGIE =================
\section{Stratégie actionnable (basée sur les données)}

\subsection{\faRocket\ Phase 1 --- Lancement (Mois 1--3)}
\begin{itemize}
  \item \textbf{Catalogue ciblé} : tapis de selle, bridons, bottes/protections, pansage, chemises anti-mouches (climat chaud : les couvertures d'hiver épaisses = stock mort), vêtements cavalier, casques
  \item \textbf{Paiements} : carte CMI + \textbf{COD} + commande WhatsApp (comportement dominant)
  \item \textbf{Boutique mobile-first}, en français + arabe + anglais
  \item \textbf{Canaux de vente} : boutiques Instagram/TikTok (62\,\% des internautes) + site marchand
\end{itemize}

\subsection{\faHandshake\ Phase 2 --- B2B (Mois 3--6)}
\begin{itemize}
  \item Cibler : 150+ clubs, 1\,000+ sorbas, entraîneurs des 7 hippodromes
  \item Offres groupées/contrats pour les clubs (sellerie, couvertures, pansage, compléments)
  \item Participer au \textbf{Salon du Cheval d'El Jadida} (150\,000 visiteurs, octobre) --- stand ou pop-up
  \item Partenariats avec les marques d'aliments déjà importées au Maroc (ex. Cavalor) en cross-sell
\end{itemize}

\subsection{\faGem\ Phase 3 --- Tourisme \& premium (Mois 6--12)}
\begin{itemize}
  \item Matériel de randonnée pour le nouvel itinéraire d'Ifrane (181 km) et les opérations plage d'Essaouira
  \item Selles/marques premium pour les clubs de polo (Marrakech, Rabat) et cavaliers pros
\end{itemize}

\subsection{\faChartLine\ Chiffres clés pour le business plan}
\begin{keybox}
\RaggedRight
\textbf{MAM (marché adressable) :} 170\,000 chevaux $\times$ $\sim$1\,000 MAD/an d'équipement = \textbf{170~M~MAD/an}\\[1mm]
\textbf{MSC (marché servable) :} 20\,000 cavaliers + 150 clubs + 1\,000 sorbas + propriétaires de chevaux de course\\[1mm]
\textbf{Croissance :} e-commerce +30\,\%/an ; filière équine croît plus vite que le PIB ; naissances +30\,\% depuis 2011
\end{keybox}

\begin{warnbox}
\begin{itemize}
  \item Droits de douane 25--30\,\% si achat hors UE $\rightarrow$ \textbf{s'approvisionner uniquement en UE (EUR.1 = 0\,\%)}
  \item COD = friction de trésorerie et risque de retours $\rightarrow$ plafonner le COD sur les gros articles, favoriser la carte CMI (47\,\%)
  \item Délais d'importation (ports de Casablanca/Tanger Med) $\rightarrow$ stock tampon, transitaire agréé
  \item Saisonnalité : pics de vente en octobre (Salon) et autour des moussems $\rightarrow$ planifier les promotions
  \item Sensibilité aux prix : panier moyen 380 MAD $\rightarrow$ points d'entrée de gamme + offre premium pour les pros
\end{itemize}
\end{warnbox}

% ================= SOURCES =================
\section*{Sources}
{\footnotesize\color{graytext}
\begin{multicols}{2}
\raggedcolumns
\begin{itemize}
  \item SOREC (sorec.ma) : élevage, courses, infrastructures
  \item Le360 (oct. 2024) : filière équine 0,6\,\% du PIB, 30\,000 emplois, contrat-programme 1 milliard MAD
  \item Oxford Business Group : croissance des courses, spectateurs Tbourida
  \item H24info / Express TV / Le Matin / Aujourd'hui le Maroc (oct. 2025) : 16\textsuperscript{e} Salon du Cheval
  \item UNESCO (ich.unesco.org) : dossier Tbourida (environ 1\,000 troupes, 330 officielles)
  \item FRMSE / Barlaman Today : trophées Tbourida, 5\,900 chevaux
  \item MesLoisirs.ma (2026) : 150+ centres, prix réels des cours et pensions
  \item FFE / Greenval (mai 2026) : itinéraires équestres d'Ifrane 181 km
  \item Observatoire Wafir E-commerce (2026) / CMI : paiements, acheteurs, panier moyen
  \item Research \& Markets (nov. 2025) : marché B2C 3,17 $\rightarrow$ 3,51 Mds \$US
  \item douane.gov.ma / customsclearance.ma / creationsocietecasablanca.com : calculs d'importation
  \item Decathlon.ma, marechalpro.com, cavalormaroc.com : prix concurrents
\end{itemize}
\end{multicols}}

\end{document}